\chapter{Symmetries, degeneracies and representations}\label{ch:intro}

\begin{watch}{\lec{1}}
Why group theory \ts{1}{51}{0:51} $\cdot$
continuous vs discrete symmetries \ts{1}{193}{3:13} $\cdot$
methane: what group theory predicts \ts{1}{824}{13:44} $\cdot$
the idea of a representation \ts{1}{1394}{23:14} $\cdot$
Wigner's theorem \ts{1}{2278}{37:58} $\cdot$
the group $\Cv{3}$ and the vector representation \ts{1}{2528}{42:08} $\cdot$
definition of a representation \ts{1}{2784}{46:24} $\cdot$
reducible vs irreducible \ts{1}{4302}{71:42} $\cdot$
equivalent representations \ts{1}{4990}{83:10}
\end{watch}

\section{Why group theory?}

Group theory is the systematic way of extracting \emph{all} the consequences of the
symmetry of a physical system, reliably and efficiently. Its computations are simple
(mostly sums and products of integers) but the concepts need care.

In condensed matter we deal almost exclusively with \emph{space-time} symmetries:
rotations, reflections, inversion, lattice translations and time reversal. Internal
symmetries ($\U{1}$, $\SU{2}$ isospin, $\SU{3}$ colour) belong mostly to high-energy
physics. A second classification is more important for the method:

\begin{description}
\item[Continuous symmetries] such as $\SO{3}$ are Lie groups. A finite rotation is
the exponential of an infinitesimal one,
\begin{equation}\label{eq:rotexp}
  \hat{U}(\vec{\theta}) = \exp\!\Bigl(-\frac{\ii}{\hbar}\,\vec{\theta}\cdot\hat{\vec{L}}\Bigr),
\end{equation}
so invariance of the Hamiltonian, $\hat{U}\hat{H}\hat{U}\herm=\hat{H}$ for all
angles, implies at first order $[\hat{H},\hat{\vec{L}}]=0$: a \emph{conservation law}.
Degeneracies and selection rules then follow from the conservation law and the
algebra of the generators, without any group theory. The $(2l+1)$-fold degeneracy
of the levels of a central potential and the dipole selection rule
$\Delta l=\pm1$ are derived this way in every quantum-mechanics course.
\item[Discrete symmetries] such as the finite set of rotations and reflections that
leave a molecule or a crystal invariant have no infinitesimal elements, hence no
generators and no conservation laws. There is no Noether theorem for them.
To go from symmetry to degeneracies and selection rules we must jump directly,
and the bridge is the representation theory of finite groups: characters, their
orthogonality, projection operators.
\end{description}

\begin{center}
\begin{tikzpicture}[node distance=1.6cm,>=Stealth,every node/.style={font=\small}]
\node (s1) {continuous symmetry};
\node[below=0.9cm of s1] (c1) {conservation law};
\node[below=0.9cm of c1] (d1) {degeneracies, selection rules};
\draw[->] (s1)--node[right]{Noether} (c1);
\draw[->] (c1)--node[right,font=\footnotesize]{algebra of generators} (d1);
\node[right=3.4cm of s1] (s2) {discrete symmetry};
\node[below=2.3cm of s2] (d2) {degeneracies, selection rules};
\draw[->,thick,blue!60!black] (s2)--node[right,text width=2.4cm,font=\footnotesize]{representation theory of finite groups} (d2);
\end{tikzpicture}
\end{center}

In this course we never need Lie-algebra techniques such as roots, Dynkin diagrams
or Young tableaux. From the rotation group we only need its irreducible representations,
which we obtain elementarily in \cref{ch:spin}. All the groups we meet are discrete
subgroups of $\Og{3}$ (rotations, reflections, inversion), lattice translations, and
time reversal.

\section{What group theory predicts: the vibrations of methane}\label{sec:methane}

The methane molecule $\mathrm{CH_4}$ has its four hydrogen atoms at the vertices of a
regular tetrahedron with the carbon atom at the centre. Its symmetry group, the
tetrahedral group $\Td$, has 24 elements.

Without symmetry considerations we can only say the following. Each of the five
atoms can move in three directions, so a distortion of the molecule is a vector with
$15$ components. In the harmonic approximation the frequencies are obtained by
diagonalising a real symmetric $15\times15$ matrix, which gives at most $15$ distinct
frequencies. Six of them vanish: three uniform translations and three rigid
rotations do not stretch any bond (these are the \emph{zero modes} of a non-linear
molecule). So at most nine non-zero frequencies remain.

A group-theory computation, which after \cref{ch:building} will take you five minutes
with pencil and paper, says much more:

\begin{keyresult}[What group theory says about $\mathrm{CH_4}$ and $\mathrm{CH_3D}$]
\begin{itemize}
\item $\mathrm{CH_4}$ has exactly \textbf{four} distinct vibrational frequencies,
      with degeneracies $1$, $2$, $3$ and $3$ ($1+2+3+3=9$).
\item All four are Raman active; only the two triply degenerate ones are infrared
      active; the molecule has no microwave (pure rotational) spectrum.
\item In deuterated methane $\mathrm{CH_3D}$ each triply degenerate frequency splits
      into one non-degenerate and one doubly degenerate frequency: six distinct
      frequencies, three non-degenerate and three doubly degenerate, all of them
      Raman and infrared active; $\mathrm{CH_3D}$ is also microwave active.
\end{itemize}
\end{keyresult}

The result of the computation is written as
\begin{equation}\label{eq:CH4result}
\begin{aligned}
  \GM  &= \irr{A}_1\oplus\irr{E}\oplus\irr{T}_1\oplus 3\,\irr{T}_2 &&(\dim 15),\\
  \GV  &= \irr{T}_2 &&(\dim 3),\qquad \GR = \irr{T}_1 \quad(\dim 3),\\
  \Gvib &= \GM\ominus\GV\ominus\GR = \irr{A}_1\oplus\irr{E}\oplus 2\,\irr{T}_2 &&(\dim 9).
\end{aligned}
\end{equation}
Every symbol here is a \emph{representation}: a rule telling how some kind of object
transforms under the symmetry operations. Loosely:
\begin{itemize}
\item $\GV$, the \emph{vector representation}, tells how ordinary (polar) vectors
      transform: a set of $3\times3$ matrices, one per symmetry operation.
\item $\GR$, the \emph{rotational} or \emph{axial-vector representation}, tells how
      pseudovectors (angular momentum, magnetic field) transform. It coincides with
      $\GV$ on proper rotations and differs by a sign on improper operations
      (reflections, inversion), because $\vec{L}=\vec{r}\times\vec{p}$ does not change
      sign under inversion while $\vec{r}$ and $\vec{p}$ do.
\item $\GM$, the \emph{mechanical representation}, tells how an arbitrary distortion
      of the molecule (a 15-component vector) transforms: $15\times15$ matrices.
\item $\Gvib$, the \emph{vibrational representation}, acts on the 9-dimensional
      space of pure vibrations, obtained by removing uniform translations (which
      transform as vectors) and rigid rotations (which transform as axial vectors).
\end{itemize}
The symbols $\irr{A}_1$, $\irr{A}_2$, $\irr{E}$, $\irr{T}_1$, $\irr{T}_2$ are the five
\emph{irreducible representations} (irreps) of $\Td$. Unlike $\GV$ or $\GM$ they do not
describe any particular object; they are abstract building blocks. Their dimension is
read off the letter: $\irr{A}$ and $\irr{B}$ are one-dimensional, $\irr{E}$ is
two-dimensional, $\irr{T}$ is three-dimensional.

\begin{keyresult}[How to read a decomposition (Wigner's theorem)]
If the representation carried by the eigenvectors of a symmetric operator (here
the dynamical matrix, later the Hamiltonian) decomposes as
$\Gamma=\bigoplus_\mu m_\mu\,\Gamma^{(\mu)}$, then
\begin{itemize}
\item the \textbf{multiplicity} $m_\mu$ is the number of distinct eigenvalues
      associated with the irrep $\Gamma^{(\mu)}$;
\item the \textbf{dimension} $d_\mu$ of the irrep is the degeneracy of each of them.
\end{itemize}
\end{keyresult}
For $\mathrm{CH_4}$, $\Gvib=\irr{A}_1\oplus\irr{E}\oplus2\,\irr{T}_2$ gives one
non-degenerate, one doubly degenerate and two triply degenerate frequencies. For
$\mathrm{CH_3D}$ the symmetry group is the smaller group $\Cv{3}$ and
$\Gvib=3\,\irr{A}_1\oplus3\,\irr{E}$ (\cref{sec:mechanical}). We prove Wigner's
theorem in an explicit example in \cref{sec:wigner}; the general proof rests on
Schur's lemma (\cref{sec:schur}).

\begin{remark}
Group theory tells you which frequencies must be degenerate, how many distinct ones
there are, and which matrix elements must vanish. It never tells you the numerical
values of the frequencies: those need physical input (force constants).
\end{remark}

\section{Groups}\label{sec:groups}

\begin{definition}
A \emph{group} is a set $G$ with a composition law $\circ$ such that
\begin{enumerate}
\item $g_1\circ g_2\in G$ for all $g_1,g_2\in G$ (closure);
\item $(g_1\circ g_2)\circ g_3=g_1\circ(g_2\circ g_3)$ (associativity);
\item there is a neutral element $e$ with $e\circ g=g\circ e=g$;
\item every $g$ has an inverse $g^{-1}$ with $g\circ g^{-1}=g^{-1}\circ g=e$.
\end{enumerate}
The number of elements $\abs{G}$ is the \emph{order} of the group.
\end{definition}
For a group of symmetry operations, $g_1\circ g_2$ means ``apply $g_2$ first, then
$g_1$''. The composition law can be anything: for the translation group it is vector
addition, $\vec{t}_1\circ\vec{t}_2=\vec{t}_1+\vec{t}_2$. Symmetry operations are
denoted as follows (Schoenflies notation):
\begin{itemize}
\item $E$: identity (do nothing). Not to be confused with the irrep $\irr{E}$.
\item $C_n$: rotation by $2\cpi/n$; $C_n^{+}$ counterclockwise, $C_n^{-}$ clockwise
      (seen from the tip of the axis).
\item $\sigma$: reflection in a plane; $\sv$ vertical (containing the main axis),
      $\sh$ horizontal, $\sd$ diagonal.
\item $I$: inversion, $\vec{r}\mapsto-\vec{r}$.
\item $S_n$: rotoreflection, a rotation $C_n$ followed by reflection in the plane
      perpendicular to the axis (\cref{sec:pointops}).
\end{itemize}

\subsection{The group \texorpdfstring{$\Cv{3}$}{C3v}}

Replacing one hydrogen of methane by deuterium leaves a triangular pyramid with the
deuterium on top. Its symmetry group $\Cv{3}$ (crystallographic symbol $3m$) has six
elements (\cref{fig:c3v}): the identity $E$, the rotations $C_3^{\pm}$ by
$\pm120^\circ$ about the vertical $z$ axis, and three vertical mirror planes
$\sigma_{\mathrm{d}1},\sigma_{\mathrm{d}2},\sigma_{\mathrm{d}3}$, each containing the
$z$ axis and one hydrogen atom. We choose the $x$ axis inside
$\sigma_{\mathrm{d}1}$.

\begin{figure}[ht]
\centering
\begin{tikzpicture}[scale=2.0,>=Stealth]
  \foreach \a/\n in {0/1,120/2,240/3}{
    \draw[dashed,gray] (\a:-1.45)--(\a:1.45);
    \node[font=\small] at (\a:1.72) {$\sigma_{\mathrm{d}\n}$};
    \fill[white,draw=black] (\a:1) circle (0.09);
    \node[font=\scriptsize] at ($(\a:1)+(\a+90:0.2)$) {H$_\n$};
  }
  \draw[thick] (0:1)--(120:1)--(240:1)--cycle;
  \fill (0,0) circle (0.06);
  \node[font=\scriptsize,below right] at (0.02,-0.02) {C, D};
  \draw[->] (0,0)--(0.7,0) node[below right,font=\small]{$x$};
  \draw[->] (0,0)--(0,0.7) node[above,font=\small]{$y$};
  \draw[->,blue!60!black] (200:0.4) arc (200:280:0.4);
  \node[blue!60!black,font=\small] at (240:0.6) {$C_3^+$};
\end{tikzpicture}
\caption{The group $\Cv{3}$ seen from above ($z$ out of the page). The deuterium
atom sits above the carbon on the $z$ axis; the three hydrogens lie below. The
mirror $\sigma_{\mathrm{d}i}$ contains the $z$ axis and hydrogen H$_i$.}
\label{fig:c3v}
\end{figure}

Its multiplication table (row $g$, column $h$, entry $g h$, i.e.\ $h$ applied first)
is
\begin{equation}\label{eq:c3vtable}
\begin{array}{c|cccccc}
 & E & C_3^+ & C_3^- & \sigma_{\mathrm{d}1} & \sigma_{\mathrm{d}2} & \sigma_{\mathrm{d}3}\\\hline
E & E & C_3^+ & C_3^- & \sigma_{\mathrm{d}1} & \sigma_{\mathrm{d}2} & \sigma_{\mathrm{d}3}\\
C_3^+ & C_3^+ & C_3^- & E & \sigma_{\mathrm{d}3} & \sigma_{\mathrm{d}1} & \sigma_{\mathrm{d}2}\\
C_3^- & C_3^- & E & C_3^+ & \sigma_{\mathrm{d}2} & \sigma_{\mathrm{d}3} & \sigma_{\mathrm{d}1}\\
\sigma_{\mathrm{d}1} & \sigma_{\mathrm{d}1} & \sigma_{\mathrm{d}2} & \sigma_{\mathrm{d}3} & E & C_3^+ & C_3^-\\
\sigma_{\mathrm{d}2} & \sigma_{\mathrm{d}2} & \sigma_{\mathrm{d}3} & \sigma_{\mathrm{d}1} & C_3^- & E & C_3^+\\
\sigma_{\mathrm{d}3} & \sigma_{\mathrm{d}3} & \sigma_{\mathrm{d}1} & \sigma_{\mathrm{d}2} & C_3^+ & C_3^- & E
\end{array}
\end{equation}
The group is \emph{not} commutative: $C_3^+\sigma_{\mathrm{d}1}=\sigma_{\mathrm{d}3}$
but $\sigma_{\mathrm{d}1}C_3^+=\sigma_{\mathrm{d}2}$. Each row and each column
contains every element exactly once (the \emph{rearrangement theorem}).

\section{Subgroups, cosets and Lagrange's theorem\beyondmark}\label{sec:cosets}

The lectures skip the structural part of group theory. Space groups, little groups and
induced representations use it constantly, so the starred sections supply it where it
first becomes useful.

A subset $H\subseteq G$ that is itself a group under the same law is a
\emph{subgroup}. For a finite subset it suffices to check closure.

\begin{definition}
For $g\in G$ the \emph{left coset} of $H$ is $gH=\{gh:h\in H\}$; the right coset is
$Hg$.
\end{definition}
Two left cosets are either identical or disjoint ($g_1H\cap g_2H\neq\emptyset$ implies
$g_2^{-1}g_1\in H$, hence $g_1H=g_2H$), and each contains $\abs{H}$ elements. So $G$ is a
disjoint union of cosets,
\[
G=g_1H\cup g_2H\cup\dots\cup g_nH,\qquad n=\abs{G}/\abs{H}.
\]
\begin{theorem}[Lagrange]
The order of a subgroup divides the order of the group. The integer
$[G:H]=\abs{G}/\abs{H}$ is the \emph{index} of $H$.
\end{theorem}
Consequences: the order of every element (the smallest $n$ with $g^n=e$) divides
$\abs{G}$; a group of prime order is cyclic. $\Cv{3}$ ($\abs{G}=6$) can only have
subgroups of order $1,2,3,6$: $\{E\}$, the three $\{E,\sigma_{\mathrm{d}i}\}$, $C_3$, and
$\Cv{3}$ itself.

\section{Representations}\label{sec:reps}

\subsection{The vector representation of \texorpdfstring{$\Cv{3}$}{C3v}}

A rotation or reflection acts linearly on vectors, so it is described by a
$3\times3$ matrix acting on the column of components $(v_x,v_y,v_z)\T$. With the axes
of \cref{fig:c3v}:
\begin{equation}\label{eq:VC3v}
\begin{gathered}
\mat{V}(E)=\one,\qquad
\mat{V}(\sigma_{\mathrm{d}1})=\begin{pmatrix}1&0&0\\0&-1&0\\0&0&1\end{pmatrix},\\[4pt]
\mat{V}(C_3^{+})=\begin{pmatrix}-\frac12&-\frac{\sqrt3}{2}&0\\ \frac{\sqrt3}{2}&-\frac12&0\\0&0&1\end{pmatrix},\qquad
\mat{V}(C_3^{-})=\begin{pmatrix}-\frac12&\frac{\sqrt3}{2}&0\\ -\frac{\sqrt3}{2}&-\frac12&0\\0&0&1\end{pmatrix},\\[4pt]
\mat{V}(\sigma_{\mathrm{d}2})=\begin{pmatrix}-\frac12&-\frac{\sqrt3}{2}&0\\ -\frac{\sqrt3}{2}&\frac12&0\\0&0&1\end{pmatrix},\qquad
\mat{V}(\sigma_{\mathrm{d}3})=\begin{pmatrix}-\frac12&\frac{\sqrt3}{2}&0\\ \frac{\sqrt3}{2}&\frac12&0\\0&0&1\end{pmatrix}.
\end{gathered}
\end{equation}
The mirror $\sigma_{\mathrm{d}1}$ contains the $x$ and $z$ axes, so it only reverses
$v_y$. How to construct the other matrices without memorising formulas is explained
in \cref{sec:active}. The matrices reproduce the multiplication table; for example
$\mat{V}(C_3^+)\mat{V}(C_3^+)=\mat{V}(C_3^-)$, because two rotations by $120^\circ$
make a rotation by $240^\circ=-120^\circ$.

\subsection{Definition}

\begin{definition}
A (linear) \emph{representation} $\Gamma$ of a group $G$ on a vector space $L$
assigns to each $g\in G$ an invertible linear operator $\mat{D}(g)$ on $L$ such that
\begin{equation}\label{eq:homom}
  \mat{D}(g_1)\,\mat{D}(g_2)=\mat{D}(g_1 g_2)\qquad\text{for all } g_1,g_2\in G.
\end{equation}
The dimension of $L$ is the \emph{dimension} of the representation. Choosing a basis
of $L$, each $\mat{D}(g)$ becomes a matrix.
\end{definition}
Whatever the composition law of the group (vector addition for translations), the
operators of a representation are always composed by \emph{matrix multiplication}.
From \eqref{eq:homom} it follows that $\mat{D}(e)=\one$ and
$\mat{D}(g^{-1})=\mat{D}(g)^{-1}$.

\begin{itemize}
\item The assignment $\mat{D}(g)=1$ (a $1\times1$ matrix) for every $g$ satisfies
      \eqref{eq:homom}. It is the \emph{identity} or \emph{trivial} representation,
      called $\irr{A}_1$ in $\Cv{3}$ and $\Td$. It is not useless: it describes
      invariant quantities.
\item A representation is \emph{faithful} if different elements get different
      matrices. The vector representation of $\Cv{3}$ is faithful; the trivial one is
      not. Faithfulness is not required.
\item There are infinitely many representations (any dimension, any basis). The
      central task of the theory is to classify them and to reduce any given one to
      a few standard pieces.
\end{itemize}

\section{Reducible and irreducible representations}\label{sec:reducible}

All six matrices \eqref{eq:VC3v} have the same block-diagonal structure: a $2\times2$
block acting on $(v_x,v_y)$ and a $1\times1$ block acting on $v_z$. Because matrix
multiplication preserves this structure, each set of blocks is itself a
representation. The $1\times1$ blocks are all equal to $1$: that is $\irr{A}_1$. The
$2\times2$ blocks form a two-dimensional representation called $\irr{E}$. We write
\begin{equation}
  \GV=\irr{A}_1\oplus\irr{E}.
\end{equation}
The symbol $\oplus$ (\emph{direct sum}) does not add matrices of different sizes; it
means that in a suitable basis all matrices take block-diagonal form with the given
blocks.

Geometrically: the $z$ axis and the $xy$ plane are \emph{invariant subspaces}. No
operation of the group takes a vector on the $z$ axis off the axis, nor a vector in
the plane out of the plane.

\begin{definition}
A representation on $L$ is \emph{reducible} if $L$ has a proper invariant subspace
($\neq\{0\},L$) and \emph{irreducible} otherwise.
\end{definition}
The representation $\irr{E}$ is irreducible: every line in the $xy$ plane is moved by
some $C_3$ rotation or some mirror.

For finite groups every representation is equivalent to a unitary one
(\cref{sec:unitarity}), and for unitary representations the orthogonal complement of an
invariant subspace is also invariant. Hence a reducible representation always
splits completely into a direct sum of irreducible ones. This is why irreps are
\emph{building blocks}.

\section{Equivalent representations}\label{sec:equiv}

The block structure depends on the basis. Rotate the $y$ and $z$ axes by $45^\circ$
about $x$, i.e.\ change basis with
\[
\mat{A}=\begin{pmatrix}1&0&0\\0&\tfrac{1}{\sqrt2}&-\tfrac{1}{\sqrt2}\\0&\tfrac{1}{\sqrt2}&\tfrac{1}{\sqrt2}\end{pmatrix},
\qquad \mat{W}(g)=\mat{A}\,\mat{V}(g)\,\mat{A}^{-1}.
\]
Then, for instance,
\[
\mat{W}(C_3^+)=\begin{pmatrix}-\frac12&-\frac{\sqrt6}{4}&-\frac{\sqrt6}{4}\\[2pt]
\frac{\sqrt6}{4}&\frac14&-\frac34\\[2pt]\frac{\sqrt6}{4}&-\frac34&\frac14\end{pmatrix},
\qquad
\mat{W}(\sigma_{\mathrm{d}1})=\begin{pmatrix}1&0&0\\0&0&-1\\0&-1&0\end{pmatrix},
\]
and no block structure is visible, although $\mat{W}$ describes exactly the same
physics in rotated axes.

\begin{definition}
Two representations $\mat{D}_1$ and $\mat{D}_2$ of $G$ are \emph{equivalent} if there
is an invertible matrix $\mat{A}$ such that
$\mat{D}_2(g)=\mat{A}\,\mat{D}_1(g)\,\mat{A}^{-1}$ for all $g\in G$.
\end{definition}
Equivalent representations are the same representation written in different bases.
This raises two practical questions: given a set of matrices, how do we know whether
it is reducible, and how do we find its irreducible components, without trying
infinitely many changes of basis? The answer is the theory of characters
(\cref{ch:building}). It also shows that a finite group has only finitely many
inequivalent irreps, as many as it has conjugacy classes: $\Cv{3}$ has three and
$\Td$ has five.

\section{Unitarity and complete reducibility\beyondmark}\label{sec:unitarity}

The lectures state the main theorems of representation theory without proof. The
starred sections supply the proofs, starting here with the two facts used in
\cref{sec:reducible}. $G$ is a finite group throughout; sums over $g$ run over all
elements.

\begin{theorem}
Every representation of a finite group is equivalent to a unitary one.
\end{theorem}
\begin{proof}
Given $\mat{D}(g)$ on $\mathbb{C}^n$ with the standard inner product $(x,y)=x\herm y$,
define the averaged inner product
\[\{x,y\}=\sum_g\bigl(\mat{D}(g)x,\mat{D}(g)y\bigr).\] It is positive
definite and invariant: $\{\mat{D}(h)x,\mat{D}(h)y\}=\sum_g(\mat{D}(gh)x,\mat{D}(gh)y)=\{x,y\}$
by the rearrangement theorem. In a basis orthonormal for $\{\cdot,\cdot\}$ the matrices
are unitary.
\end{proof}

\begin{theorem}[Maschke]
If $W$ is an invariant subspace of a unitary representation, so is its orthogonal
complement $W^\perp$. Hence every representation is a direct sum of irreps.
\end{theorem}
\begin{proof}
If $y\in W^\perp$ and $x\in W$, then $(x,\mat{D}(g)y)=(\mat{D}(g)^{-1}x,y)=0$, because
$\mat{D}(g)^{-1}x=\mat{D}(g^{-1})x\in W$. Iterate on $W$ and $W^\perp$ until the pieces are
irreducible (the dimension decreases at each step).
\end{proof}

\exercisesection

\begin{exercise}\label{ex:1:table}
Verify three entries of the multiplication table \eqref{eq:c3vtable} by multiplying
the matrices \eqref{eq:VC3v}: $C_3^+\sigma_{\mathrm{d}1}$,
$\sigma_{\mathrm{d}1}C_3^+$ and $\sigma_{\mathrm{d}2}\sigma_{\mathrm{d}1}$. Interpret
the last one geometrically.
\end{exercise}

\begin{solution}
With the matrices \eqref{eq:VC3v} (only the upper-left $2\times2$ blocks matter; the
$zz$ entry is always $1$):
\[
\mat{V}(C_3^+)\mat{V}(\sigma_{\mathrm{d}1})=
\begin{pmatrix}-\frac12&-\frac{\sqrt3}{2}\\ \frac{\sqrt3}{2}&-\frac12\end{pmatrix}
\begin{pmatrix}1&0\\0&-1\end{pmatrix}=
\begin{pmatrix}-\frac12&\frac{\sqrt3}{2}\\ \frac{\sqrt3}{2}&\frac12\end{pmatrix}
=\mat{V}(\sigma_{\mathrm{d}3}),
\]
\[
\mat{V}(\sigma_{\mathrm{d}1})\mat{V}(C_3^+)=
\begin{pmatrix}-\frac12&-\frac{\sqrt3}{2}\\ -\frac{\sqrt3}{2}&\frac12\end{pmatrix}
=\mat{V}(\sigma_{\mathrm{d}2}),
\]
\[
\mat{V}(\sigma_{\mathrm{d}2})\mat{V}(\sigma_{\mathrm{d}1})=
\begin{pmatrix}-\frac12&\frac{\sqrt3}{2}\\ -\frac{\sqrt3}{2}&-\frac12\end{pmatrix}
=\mat{V}(C_3^-).
\]
Geometrically, reflecting first in a line and then in a second line making the angle
$\alpha$ with it is a rotation by $2\alpha$. From $\sigma_{\mathrm{d}1}$ ($0^\circ$) to
$\sigma_{\mathrm{d}2}$ ($120^\circ$) the angle is $120^\circ$, so the product is a rotation
by $240^\circ=-120^\circ$, i.e.\ $C_3^-$.
\end{solution}

\begin{exercise}\label{ex:1:subgroups}
List all subgroups of $\Cv{3}$. Which of them are commutative?
\end{exercise}

\begin{solution}
By Lagrange's theorem the orders can be $1,2,3,6$: $\{E\}$; the three groups
$\{E,\sigma_{\mathrm{d}i}\}$; $C_3=\{E,C_3^+,C_3^-\}$; and $\Cv{3}$. All proper subgroups have
prime or unit order, hence are cyclic and commutative; $\Cv{3}$ itself is not
commutative.
\end{solution}

\begin{exercise}\label{ex:1:A2}
Show that assigning $+1$ to $E,C_3^\pm$ and $-1$ to the three mirrors is a
representation of $\Cv{3}$ (called $\irr{A}_2$). Is it faithful? Give a physical
quantity that transforms according to $\irr{A}_2$.
\end{exercise}

\begin{solution}
The assignment is $g\mapsto\det\mat{V}(g)$, and $\det$ is multiplicative:
$\det(\mat{V}(g_1)\mat{V}(g_2))=\det\mat{V}(g_1)\det\mat{V}(g_2)$. (Equivalently: rotation
times rotation and mirror times mirror are rotations, rotation times mirror is a mirror,
as the table \eqref{eq:c3vtable} confirms.) It is not faithful: its kernel is $C_3$.
Quantities transforming as $\irr{A}_2$: the $z$ component of an axial vector ($R_z$,
$L_z$, $B_z$), which is invariant under the rotations about $z$ and reversed by vertical
mirrors.
\end{solution}

\begin{exercise}\label{ex:1:count}
A planar molecule $\mathrm{BF_3}$ has four atoms. How many vibrational degrees of
freedom does it have? What about the linear molecule $\mathrm{CO_2}$? (For a linear
molecule, rotation about the molecular axis does not move the atoms.)
\end{exercise}

\begin{solution}
$\mathrm{BF_3}$: $3\times4=12$ degrees of freedom, minus 3 translations and 3 rotations:
$6$ vibrational. $\mathrm{CO_2}$ (linear): $9-3-2=4$.
\end{solution}

\begin{exercise}\label{ex:1:equiv}
Show that the matrices $\mat{W}(g)$ of \cref{sec:equiv} have the same traces as the
matrices $\mat{V}(g)$. Then prove that $\tr(\mat{A}\mat{B}\mat{A}^{-1})=\tr\mat{B}$
for any invertible $\mat{A}$.
\end{exercise}

\begin{solution}
$\tr\mat{W}(C_3^+)=-\tfrac12+\tfrac14+\tfrac14=0=\tr\mat{V}(C_3^+)$ and
$\tr\mat{W}(\sigma_{\mathrm{d}1})=1+0+0=1=\tr\mat{V}(\sigma_{\mathrm{d}1})$. In general
$\tr(\mat{A}\mat{B}\mat{A}^{-1})=\tr(\mat{A}^{-1}\mat{A}\mat{B})=\tr\mat{B}$ by the cyclic
property $\tr(\mat{X}\mat{Y})=\sum_{ij}X_{ij}Y_{ji}=\tr(\mat{Y}\mat{X})$.
\end{solution}
